\section{Full model, measurable access, and minimax quantifiers}
\label{app:model}

\subsection{Measurable experiment and statistical target}

Fix once and for all a nonempty infinite standard Borel measurable space
\((\mathcal X,\mathcal A)\).  This fixed nonempty infinite standard Borel measurable space is shared by every public pair and law below.  Equip
\[
 \mathcal Z=\mathcal X\times\mathbb R\times\mathbb R
\]
with the product sigma-field
\(\mathcal A\otimes\mathcal B(\mathbb R)\otimes
\mathcal B(\mathbb R)\).  The coordinates
\(X:\mathcal Z\to\mathcal X\), \(T,Y:\mathcal Z\to\mathbb R\) are
measurable, and all conditional expectations are ordinary versions on this
probability space.  Public represented functions are jointly supplied with
measurable evaluation maps \(x\mapsto g_j(x)\).  This standard-Borel
condition is part of the theorem field before \(n\), the budget tuple, the
public pair, and the compatible law are chosen.  It excludes pathological
nonstandard measurable spaces and finite ambient spaces from the matching
I-stratum theorem, while including countably infinite and uncountable
standard Borel spaces used in ordinary ML models.

For a fixed regularity tuple
\[
 \rho=(B_0,B_{\mathcal G},v_-,v_+,c_{\rm rad}),\qquad
 B_0\ge3,\quad B_{\mathcal G}>0,\quad
 0<v_-\le v_+,\quad c_{\rm rad}>0,
\]
a law in \(\cP(\Gamma)\) satisfies \eqref{eq:plm},
\(\E(u\mid X)=0\), \(\E(\varepsilon\mid X,T)=0\), and
\(v:=\E u^2\in[v_-,v_+]\).  The conditional variance
\(\operatorname{Var}(T\mid X)=\E(u^2\mid X)\) may vary with \(X\)
and may vanish on subsets of the covariate space.  It also satisfies
\[
 |Y|,|T|,|\beta_0|,|\mu_0(X)|,|\pi_0(X)|,|u|,|\varepsilon|
 \le B_0,\qquad \beta_0\in[-3,3].
\]
The target is the statistical coefficient \(\beta_0\) in this conditional
mean model.  Causal interpretation would require additional assumptions and
is not asserted.

Every admissible public pair also satisfies the fixed pointwise candidate
envelope
\[
 \sup_{g\in\cG_\mu\cup\cG_\pi}|g(x)|\le B_{\mathcal G}
 \qquad\text{for every }x\in\mathcal X.
 \tag{A.0}
\]
Every function in the supplied countable representations obeys the same
bound.  Thus each public difference class has envelope \(2B_{\mathcal G}\).
This condition is part of the fixed theorem field before the outer public-pair
supremum; it is not inferred from approximation or localized complexity.

\subsection{Positive-radius public complexity and M1 access}

For a public class \(\cG\), let
\(\Delta\cG=\{g-g':g,g'\in\cG\}\).  For any class \(\mathcal H\) and
\(r\ge0\), put
\[
 Z_{\mathcal H}(r):=
   \sup_{\substack{h\in\mathcal H\\ \norm h_{L_2(P_X)}\le r}}
   \abs{\frac1n\sum_{i=1}^n\epsilon_i h(X_i)},\qquad
 \mathfrak R^{*}_{n,P_X}(r;\mathcal H)
 :=\E^{*}_{X_{1:n},\epsilon_{1:n}} Z_{\mathcal H}(r).
\]
Here \(X_{1:n}\sim P_X^n\), the \(\epsilon_i\)'s are independent Rademacher
signs independent of the sample, and the outer expectation \(\E^*Z\) is the
infimum of \(\E U\) over all measurable majorants \(U\ge Z\).  When
\(Z_{\mathcal H}(r)\) is measurable, outer and ordinary expectation agree.
If \(\cG^{\rm rep}=\{g_j:j\in\mathbb N\}\), its countable represented difference core is
\(\Delta\cG^{\rm rep}=\{g_j-g_k:j,k\in\mathbb N\}\).  Its supremum is an ordinary measurable random variable, and we write its ordinary-expectation
complexity as \(\mathfrak R_{n,P_X}(r;\Delta\cG^{\rm rep})\).  Pointwise
set inclusion and monotonicity of outer expectation give
\[
 \mathfrak R_{n,P_X}(r;\Delta\cG^{\rm rep})
 \le \mathfrak R^{*}_{n,P_X}(r;\Delta\cG).
\]
Equality is not required, and sequential M1 alone does not provide it.  At
\(r=0\) the unquotiented quantities above remain defined.  No
star-shapedness or scaling closure is assumed or used.  The stochastic
constraints are exactly
\begin{align}
 \sup_{r\ge s,\ r>0}
 \frac{\mathfrak R^{*}_{n,P_X}(r;\Delta\cG_\mu)}{3r}
 &\le c_{\rm rad}s,\nonumber\\
 \sup_{r\ge t,\ r>0}
 \frac{\mathfrak R^{*}_{n,P_X}(r;\Delta\cG_\pi)}{3r}
 &\le c_{\rm rad}t.
 \tag{A.1}
\end{align}
No positive-radius quotient at \(r=0\) is defined or used.  When a declared
stochastic budget is zero, (A.1) directly forces the corresponding full-class
outer complexity, and hence every represented-core complexity, to vanish at
each positive radius.  Later upper-proof calls therefore use core base radius
zero directly; they never evaluate a radius-zero quotient.

M1 means that the learner receives the represented public instance
\[
 \Gamma=(\cG_\mu,\cG_\pi,\mathbf g_\mu,\mathbf g_\pi),\qquad
 \mathbf g_\nu=(g_{\nu,j})_{j\in\mathbb N},\quad \nu\in\{\mu,\pi\},
\]
where \(\mathbf g_\nu\) is an ordered countable representation of
\(\cG_\nu\).  It also receives the fixed
tuple including \(B_{\mathcal G}\), and the four declared budgets before
observing data.  The representation has measurable evaluations, obeys the
same envelope, and has the following sequential property: for every
full-class element \(g\in\cG\), there is a sequence \(g_{j_k}\) from the
representation that converges pointwise everywhere to \(g\) on \(\mathcal X\).
The common envelope and dominated convergence make the represented and
full-class population \(L_2(P_X)\) approximation infima equal; finite-sample
continuity gives the corresponding empirical fit and calibration transfers.
Empirical approximate minimizers select the first represented index within
\(n^{-12}\) of a countable infimum.
Scalar coefficients use \(\mathbb Q\cap[-3,3]\); finite-dimensional convex
objectives use objective-dense rational parameterizations.  Countable
infima, first qualifying indices, finite arithmetic, and clipping are
measurable, so every fitted rule is Borel.  M1 does not supply nuisance truths,
population best approximants, realized errors, a hard label, or an
optimization oracle.

\begin{definition}[Indexed compatible-law class]
\label{def:indexed-compatible-laws}
Fix \(\rho,n,a,s,b,t\), and let
\(\Gamma=(\cG_\mu,\cG_\pi,\mathbf g_\mu,\mathbf g_\pi)\) be a represented
public instance obeying the envelope and measurable sequential-M1 interface
above.  We call \((\cG_\mu,\cG_\pi)\) its class-pair projection.  The indexed compatible-law class
\(\mathcal P_{\rho,n,a,s,b,t}(\Gamma)\) is the set of probability laws
\(P\) on \(\mathcal Z\) for which there are
\(\beta_0\in[-3,3]\),
\(\mathcal A/\mathcal B(\mathbb R)\text{-measurable functions }\mu_0,\pi_0:\mathcal X\to\mathbb R\),
and residuals
\[
 u=T-\pi_0(X),\qquad
 \varepsilon=Y-\beta_0T-\mu_0(X),
\]
such that the PLM representation \eqref{eq:plm}, the conditional-centering
conditions
\(\E(u\mid X)=0\) and \(\E(\varepsilon\mid X,T)=0\) hold
\(P\)-almost surely, and the average residual-variance condition
\(v:=\E_P u^2\in[v_-,v_+]\) holds.  The shared law and target envelopes are
\[
 |Y|,|T|,|\beta_0|,|\mu_0(X)|,|\pi_0(X)|,|u|,|\varepsilon|
 \le B_0,
\]
and the public classes and their representations retain the pointwise
candidate envelope (A.0), measurable evaluation maps, and the
pointwise-everywhere sequential-M1 approximation property.  Under this
law's covariate marginal \(P_X\), the approximation-budget conditions are
\[
 \inf_{g\in\cG_\mu}\norm{g-\mu_0}_2\le a,
 \qquad
 \inf_{g\in\cG_\pi}\norm{g-\pi_0}_2\le b,
\]
and the localized-complexity conditions are
\[
 \sup_{r\ge s,\ r>0}
 \frac{\mathfrak R^{*}_{n,P_X}(r;\Delta\cG_\mu)}{3r}
 \le c_{\rm rad}s,
 \qquad
 \sup_{r\ge t,\ r>0}
 \frac{\mathfrak R^{*}_{n,P_X}(r;\Delta\cG_\pi)}{3r}
 \le c_{\rm rad}t.
\]
The target \(\beta(P)\) is uniquely identified: if two such conditional-mean
PLM representations describe the same law, their conditional-mean difference
is \(\gamma T+f(X)=0\) almost surely.  Conditional centering identifies
\(\pi_0(X)=\E(T\mid X)\), hence the same residual \(u\), in both
representations.  Multiplying the difference by \(u\) and taking
expectations gives \(0=\gamma\E u^2\), since
\(\E\{f(X)u\}=\E\{\pi_0(X)u\}=0\).  Because
\(\E u^2\ge v_->0\), \(\gamma=0\), so the two coefficients agree.
Apart from recording this consequence, this
definition merely collects the model, envelope, approximation,
positive-radius complexity, and public measurable-access predicates already
stated above; it changes none of them.  In the main text and the rest of the
appendices we suppress the fixed indices \(\rho,n,a,s,b,t\) and write
\(\cP(\Gamma)\).
\end{definition}

\subsection{Q1 order and allowed sample-size dependence}

For each integer \(n\ge1\) and budget tuple, define
\(\GammaSet_{\rho,n}(a,s,b,t)\) as the collection of represented public
pairs on this same fixed ambient space satisfying the candidate-envelope and
measurable sequential-M1 fields above and having
\(\mathcal P_{\rho,n,a,s,b,t}(\Gamma)\ne\varnothing\).  The ambient space is not a coordinate
of \(\Gamma\) and is not varied by the outer supremum.
Q1 is the order printed in \eqref{eq:risk}:
\[
 \sup_{\Gamma\in\GammaSet_{\rho,n}(a,s,b,t)}
 \inf_{\delta_\Gamma\ {\rm Borel}}
 \sup_{P\in\cP(\Gamma)}
 \E_P|\delta_\Gamma(Z_{1:2n})-\beta(P)|.
\]
The public pair and a lower hard family may depend on \(n\) and the declared
budgets, but they are fixed before the data and before any latent label,
coordinate, or sign is drawn.  Dependence on realized data is prohibited.
No single class pair is required to serve all sample sizes or all budget
cells.  For each fixed \(\Gamma\), the decision rule may depend on
\(\Gamma,\rho,a,s,b,t\) but not on \(P\), proof-only comparators, or
realized errors.  Hence the outer public-pair supremum is taken only after the
fixed-pair rule and its uniform law bound are established.

The total sample size is \(N=2n\), while \(q=n^{-1/2}\); replacing \(q\) by
\(N^{-1/2}\) changes only a fixed factor.  Risk is expected absolute error.
All interior constants may depend on fixed \(\rho\), including
\(B_{\mathcal G}\), but not on
\(n,a,s,b,t,\Gamma,P\).  At \(c_{\rm rad}=0\), (A.1) collapses every legal
difference class in \(L_2(P_X)\), and the positive-budget nuisance lower is
false.  Thus the theorem is pointwise for \(c_{\rm rad}>0\), with constants
that may degenerate as \(c_{\rm rad}\downarrow0\); it is not uniform at that
boundary.
